\chapter{Particle Amplitude Scales and Stability Validation}
\label{app:particle_amplitude}
\addtocontents{toc}{\protect\setcounter{tocdepth}{0}}

\noindent\textbf{Subject:} Verification of the Hydrodynamic Mass Model via Amplitude Calculation.

\section{The Objective}
In \hyperref[app:fcc_constants]{Appendix \ref*{app:fcc_constants}}, we derived the holographic vacuum condensate density ($\rho_{\mathrm{IC}} \approx 1.8074 \times 10^{53}$ kg/m$^3$) of the vacuum. In \hyperref[sec:inertial_mass]{Section \ref*{sec:inertial_mass}}, we defined Inertial Mass ($m$) as the hydrodynamic impedance of this condensate.

\noindent\textbf{Objective:} To validate this model, we must calculate the physical vibration amplitude ($r_{\mathrm{amp}}$) of known particles. If the theory is consistent:
\begin{itemize}
    \item \textbf{The Heaviest Particle (Top Quark):} Must fit within the lattice node ($r_{\mathrm{amp}} < r_{\mathrm{D}}$).
    \item \textbf{The Lightest Particle (Electron):} Must be physically resolvable above the substrate lattice size ($r_{\mathrm{amp}} > \ell_{\mathrm{S}}$).
\end{itemize}

\section{The Mass-Amplitude Relation}
We model a fundamental particle as a spherical lattice node of radius $r_{\mathrm{D}}$ undergoing radial pulsation (``Breathing Mode'') within the vacuum condensate defined in \hyperref[app:fcc_constants]{Appendix \ref*{app:fcc_constants}}.

\subsection{The Impedance Equation}
From \hyperref[sec:inertial_mass]{Section \ref*{sec:inertial_mass}}, assuming stable phase resonance ($\sin(\Delta\phi)=1$):
\begin{equation}
    m = \rho_{\mathrm{IC}} \cdot V_{\mathrm{eff}}
\end{equation}

\subsection{The Effective Volume ($V_{\mathrm{eff}}$)}
For a sphere pulsating with radial amplitude $r_{\mathrm{amp}}$, the swept volume is the Surface Area times the Amplitude:
\begin{equation}
    V_{\mathrm{eff}} \approx A_{\mathrm{surf}} \cdot r_{\mathrm{amp}} = (4\pi r_{\mathrm{D}}^2) \cdot r_{\mathrm{amp}}
\end{equation}

\subsection{The Solution for Amplitude}
By substituting the Fundamental Impedance Constant ($K_{\mathrm{D}} = \rho_{\mathrm{IC}} \cdot 4\pi r_{\mathrm{D}}^2 \approx 126.07$ kg/m) derived in \hyperref[app:fcc_constants]{Appendix \ref*{app:fcc_constants}}, we simplify the relation and solve directly for the unknown vibration scale:
\begin{equation}
    r_{\mathrm{amp}} = \frac{m}{K_{\mathrm{D}}}
\end{equation}

\section{Validation Case A: The Top Quark (Upper Limit)}
The Top Quark ($m_{\mathrm{t}} \approx 173$ GeV/$c^2$) is the heaviest known fundamental particle. It tests the structural capacity of the lattice.

\begin{itemize}
    \item \textbf{Mass:} $m_{\mathrm{t}} \approx 3.0838 \times 10^{-25}$ kg
\end{itemize}

\noindent\textbf{Calculation:}
\begin{equation}
    r_{\mathrm{amp}} = \frac{3.0838 \times 10^{-25}}{126.07} \approx 2.4461 \times 10^{-27} \text{ m}
\end{equation}

\noindent\textbf{Stability Check (Occupancy Ratio):}
\begin{equation}
    \text{Ratio} = \frac{2.4461 \times 10^{-27}}{7.4504 \times 10^{-27}} \approx 0.3283 \quad (32.83\%)
\end{equation}

\noindent\textbf{Conclusion:} The Top Quark utilizes $32.83\%$ of the available lattice space. This explicitly confirms structural stability ($<100\%$) and leaves geometrical headroom for heavier, transient intermediate vector bosons. It physically guarantees that the Standard Model mass spectrum operates strictly within the linear elastic regime of the active vacuum lattice framework.

\section{Validation Case B: The Electron (Lower Limit)}
The Electron ($m_{\mathrm{e}} \approx 0.511$ MeV/$c^2$) is the lightest charged lepton. It tests the resolution limit of the substrate.

\begin{itemize}
    \item \textbf{Mass:} $m_{\mathrm{e}} \approx 9.1094 \times 10^{-31}$ kg
\end{itemize}

\noindent\textbf{Calculation:}
\begin{equation}
    r_{\mathrm{amp}} = \frac{9.1094 \times 10^{-31}}{126.07} \approx 7.2257 \times 10^{-33} \text{ m}
\end{equation}

\noindent\textbf{Resolution Check (Substrate Resolution Ratio):}
Using the fundamental pixel size ($\ell_{\mathrm{S}} \approx 1.6163 \times 10^{-35}$ m):
\begin{equation}
    \text{Ratio} = \frac{7.2257 \times 10^{-33}}{1.6163 \times 10^{-35}} \approx 447.05
\end{equation}

\noindent\textbf{Conclusion:} The Electron vibrates with an amplitude approximately 447 times larger than the fundamental substrate lattice ($\ell_{\mathrm{S}}$). This confirms the electron is physically ``real'' and mathematically resolvable above the Planck-scale floor. It mechanically explains why the electron appears dimensionless to our macroscopic instruments ($10^{-33}$ m is substantially below the Large Hadron Collider resolution limit of $10^{-20}$ m), yet generates a finite, continuous hydrodynamic mass.

\section{Summary}
The mass hierarchy is a direct consequence of the Vibrational Amplitude allowed by the Soft Floor geometry.

\begin{table}[H]
\centering
\caption{Particle Amplitudes and Lattice Constraints}
\vspace{0.2cm}
\label{tab:particle_amplitudes}
\renewcommand{\arraystretch}{1.5}
\setlength{\tabcolsep}{4pt}
\small
\begin{tabular}{@{}llclp{0.35\textwidth}@{}}
\toprule
\textbf{Particle} & \textbf{Mass (kg)} & \textbf{Vibration ($r_{\mathrm{amp}}$)} & \textbf{Lattice Occupancy} & \textbf{Status} \\
\midrule
Top Quark & $3.0838 \times 10^{-25}$ & $2.4461 \times 10^{-27}$ m & $32.83\%$ ($0.328 r_{\mathrm{D}}$) & Valid. Operating within geometric limit. \\
Electron & $9.1094 \times 10^{-31}$ & $7.2257 \times 10^{-33}$ m & $\ll 0.001\%$ ($447 \ell_{\mathrm{S}}$) & Valid. Resolvable above substrate floor. \\
\bottomrule
\end{tabular}
\end{table}

This validates that the TFT vacuum parameters derived in \hyperref[app:fcc_constants]{Appendix \ref*{app:fcc_constants}} provide a mechanically consistent habitat for the Standard Model particles.

\vspace{0.5em}
\noindent\textbf{Note on Constant Topological Resonance vs. Emergent Mass:}\\
Inertial mass is strictly generated by an FCC node radially squeezing and expanding (a topological ``Breathing Mode'') against the interstitial condensate. While the displacement amplitude ($r_{\mathrm{amp}}$) dictates the specific hydrodynamic impedance (and therefore the macroscopic rest energy $E_0$ and emergent Proper Time clock $t_{\mathrm{p}}$) of a specific particle, the physical resonator is always the fixed underlying node.

\begin{itemize}
    \item \textbf{The Photon Limit (Spin-1):} A continuous transverse photon strictly requires a minimum wavelength equal to the node diameter ($2r_{\mathrm{D}}$) to complete a standard $360^{\circ}$ phase cycle. This kinematic limit is physically identical to the structural resonant frequency of the active node itself ($f_{\gamma} \equiv f_{\mathrm{D}} = c_{\mathrm{E}}/2r_{\mathrm{D}}$), as derived in \hyperref[sec:coherence_boundary]{Section \ref*{sec:coherence_boundary}}.
    
    \item \textbf{The Fermion Cycle (Spin-1/2):} However, localized matter operates as a standing wave. To complete one full physical breathing cycle and return to its original topological identity, the internal mechanical phase-twist must traverse the diameter outward (Expansion = $2r_{\mathrm{D}}$, $360^{\circ}$) and retract inward (Contraction = $2r_{\mathrm{D}}$, $360^{\circ}$).
\end{itemize}

This total kinematic path of $4r_{\mathrm{D}}$ means all fundamental matter shares an identical, constant internal topological resonance of:
\begin{equation}
    f_{\mathrm{m}} = \frac{c_{\mathrm{E}}}{4r_{\mathrm{D}}} \equiv \frac{1}{2} f_{\mathrm{D}}
\end{equation}

This mechanical necessity perfectly derives the $720^{\circ}$ (Spin-1/2) phase requirement of all fermions. The fermion's foundational structural frequency is exactly half the photon limit ($f_{\mathrm{m}} = \frac{1}{2} f_{\gamma}$) because its standing-wave topology requires twice the internal transit to complete a cycle. Meanwhile, the macroscopic Proper Time clock ($t_{\mathrm{p}}$) remains strictly governed by the specific amplitude's drag against the vacuum condensate.
